% Drop-in completion of the empty "Shuttling and Landau-Zener tunnelling"
% section of orswan/opticalLattices Optical-Lattices.tex, written to match the
% manuscript's notation (natural units hbar = 1, 2m = 1) and style.
% Numerics validating eq. (LZloss) below: lz_sweep.py driving lattice1D.py.

\section{Shuttling and Landau-Zener tunnelling}

We now analyse what happens when the quasimomentum sweeps through a
band-edge degeneracy too quickly for adiabaticity to hold.  The result is a
closed-form expression for the fraction of atoms lost during a shuttle,
which sets the practical limit on how fast atoms can be moved.

\subsection{Reduction to two levels}

Specialize to a minimal 1D lattice with wavenumber $k$ and depth $A$, and
consider the swept Hamiltonian of eq. (\ref{sweepH}) restricted to the
quasimomentum subspace $\mathcal{S}_q$:
$$H_{Q(t)} = \left(p - Q(t)\right)^2 + A\cos(kx), \qquad Q(t) = m\,\dot a(t).$$
(For gravitational driving, $Q(t) = \half g t$ in our units, by the results of
the previous section.)  In the momentum basis $\left|q + nk\right>$ the kinetic
term is diagonal with entries $\left(P_n - Q(t)\right)^2$, $P_n := q + nk$, and
the lattice couples neighboring momentum states with matrix element $A/2$ ---
this is precisely the tridiagonal structure noted in the momentum space section.

Tunnelling can only occur where two diabatic (lattice-off) energies cross, i.e.
where $\left(P_n - Q\right)^2 = \left(P_{n+1} - Q\right)^2$, which happens when
$Q - q$ is a half-integer multiple of $k$: the zone edge.  Near such a crossing
all other momentum states are separated by energies of order $k^2$, so provided
$A \ll k^2$ we may project onto the two crossing states.  The projected
Hamiltonian is
$$H_2(t) = \begin{bmatrix} \left(P_n - Q(t)\right)^2 & A/2 \\ A/2 &
\left(P_{n+1} - Q(t)\right)^2 \end{bmatrix},$$
and the diabatic splitting
$$\Delta E(t) = \left(P_{n+1}-Q(t)\right)^2 - \left(P_n-Q(t)\right)^2
= k\left(P_n + P_{n+1} - 2Q(t)\right)$$
passes through zero linearly, with slew rate
$$\left|\frac{d\,\Delta E}{dt}\right| = 2k\,|\dot Q|.$$
At the crossing the adiabatic eigenvalues avoid each other by the gap $A$
(twice the coupling $A/2$), which is the band gap between the two lowest bands
at the zone edge --- consistent with first-order degenerate perturbation theory
in the band structure picture, valid in the same regime $A \ll k^2$.

\subsection{The Landau-Zener formula}

A two-level system with constant coupling $V$ whose diabatic splitting is swept
linearly through zero at rate $\alpha = |d\,\Delta E/dt|$ is the classic
Landau-Zener problem, with diabatic transition probability
$$P_{LZ} = \exp\left(-\frac{2\pi V^2}{\alpha}\right).$$
(I will not rederive this here; the standard derivation proceeds via
parabolic cylinder functions.)  With $V = A/2$ and $\alpha = 2k|\dot Q|$ we
obtain the loss per zone-edge crossing:
\begin{equation}
P_{\text{tunnel}} = \exp\left(-\frac{\pi A^2}{4 k |\dot Q|}\right)
\;\xrightarrow{\;Q = \frac{1}{2}gt\;}\;
\exp\left(-\frac{\pi A^2}{2 k g}\right)
\label{LZloss}
\end{equation}
A diabatic transition here means the atom keeps its momentum while the lattice
sweeps on without it: the atom has slipped out of the lattice.  (Strictly it
is excitation to a higher band of the swept lattice; for a monotonic shuttle
this is loss because, as shown below, the left-behind branch never rejoins
the transported one.)  The
adiabatic fraction $1 - P_{\text{tunnel}}$ is carried along --- each crossing
transfers momentum $k$ (in units where $\hbar=1$) per atom, which is how a swept lattice
accelerates atoms in the first place.

Two limits.  Deep lattice or gentle sweep
($A^2 \gg k\dot Q$): losses are exponentially small and the lattice shuttles
faithfully.  Shallow lattice or hard sweep ($A^2 \ll k\dot Q$): essentially
everything tunnels, and the lattice slides through the atoms leaving them
(almost) untouched.  Note the figure of merit is $A^2/\dot Q$, so halving the
lattice depth requires quartering the acceleration to hold losses fixed.

\subsection{Numerical check}

Equation (\ref{LZloss}) is a two-level approximation to an infinite-level
problem, so it should be checked.  The script \texttt{lz\_sweep.py} drives the
momentum-space solver of \texttt{lattice1D.py} (which integrates the full
$(2n+1)$-state Schr\"odinger equation with no two-level reduction) through a
single gravity-driven zone-edge crossing, then measures the ground-band
survival by projecting the final state onto band eigenstates in the
lattice's co-moving frame.  Across
lattice depths $A = aa\,g_0$ with $aa \in [2, 23]$ and drives
$g \in \{2, 4, 8\}\,g_0$ (where $g_0 = 0.0295$ is the code's strontium
gravity unit; the dimensionless exponent $\pi A^2/2kg$ spans
$[0.012, 6.1]$) --- losses spanning $0.2\%$ to $99\%$ --- the
simulated loss matches eq. (\ref{LZloss}) with absolute deviations between
$10^{-5}$ and $10^{-2}$ (largest at the deepest lattices and hardest drives;
relative deviations reach $\sim 8\%$ at the smallest losses, where the
denominator is tiny), consistent in magnitude with
finite-sweep endpoint effects and numerical truncation (the cause was not
isolated; finite-time LZ corrections are oscillatory, so no claim is made
about their sign).  Concretely, the projection observable is the instantaneous ground band of
the swept Hamiltonian $H_{Q_f} = (p-Q_f)^2 + V$ within the fixed
quasimomentum sector; \texttt{SFrame} in \texttt{lattice1D.py} implements
exactly this by transforming to the co-moving gauge and projecting there ---
equivalent statements in two frames, and \emph{not} a projection onto the
eigenstates of the unshifted $H_0$.

\subsection{Multiple crossings and where the formula breaks}

For a sustained sweep the quasimomentum crosses a zone edge once per Bloch
period, and if the crossings are well separated the losses compound
geometrically: after $N$ crossings the surviving fraction is
$\left(1 - P_{\text{tunnel}}\right)^N$.  For a monotonic sweep this
compounding needs less than one might fear: successive crossings involve
\emph{different} momentum pairs, so amplitude left behind at one crossing
never meets the transported branch again --- the outgoing channels are
distinct, and no interference or dephasing assumption is required.  What the
compounding does require is that the crossings be well separated, so each is
an isolated Landau-Zener event; when the sweep is hard enough that the
transition windows of successive crossings overlap, the isolated-crossing
picture itself breaks down.  The full solver captures that regime; the
formula does not.  (A non-monotonic sweep --- reversals, or added couplings
that close a loop in momentum space --- can recombine amplitudes and produce
genuine St\"uckelberg interference; none is present here.)  It likewise breaks when $A \sim k^2$, where the two-level
projection fails because the gap to the second excited band is no longer
large compared to $A$.

An interesting question in the spirit of the earlier sections: given a fixed
total momentum transfer and a fixed time budget, what trajectory $a(t)$
minimizes total tunnelling loss?  Equation (\ref{LZloss}) says loss is incurred
almost entirely near the crossings, suggesting a bang-bang profile --- sweep
hard between crossings, ease through them --- with the optimal schedule set
by a trade-off between the LZ exponent at each crossing and the sprint speed
required between crossings to meet the time budget (the number of crossings
itself is fixed by the total momentum transfer).
I leave this optimization, which is directly relevant to large momentum
transfer atom optics, for future work.
